\chapter{From data to a relational database}
\label{learn:foundations}
\learning{Understand tables, rows, keys, relationships, SQL statements, client/server sessions and durable storage. Learn which part of MiniSQL performs each job before installing anything.}

\section{What a database gives you}
\index{database}\index{relational model}
Imagine keeping a small shop in a spreadsheet. One sheet lists customers,
another products, and a third orders. A product's name should be recorded in
one authoritative place, not copied into hundreds of order rows. Every order
must identify a real customer. Two people should not accidentally sell the
same last item because they both read yesterday's stock count. If a process
stops, you need a defined way to recover committed work.

A relational database organizes facts into \emph{tables} and lets you express
relationships between them. A table has named \emph{columns}; each column has
a declared type and constraints. Each \emph{row} records one instance of the
kind of thing that the table represents. A query produces a result table.
A transaction groups related changes into one decision to keep or undo the
changes. These ideas are independent of any particular interface.

MiniSQL implements this model in native programs written in MiniLang. It
includes a persistent server, a console client, a Windows Workbench,
offline maintenance utilities, and Python and Java connectors. It also
implements its own parser, planner, executor and disk storage. Its SQL is
not SQLite's SQL, and a MiniSQL database directory is not an SQLite file.
Choose tools designed for MiniSQL and read the dialect details in this book.

\section{Tables, rows and types}
\index{table}\index{row}\index{column}
The following table records three products. A product identifier answers
``which product?''; the name describes it; the stock records a quantity.

\begin{table}[htbp]\centering
\caption{A first relation: selected columns from the teaching products table}
\begin{tabular}{r l r}\toprule
product\_id & name & stock\\\midrule
101 & Mechanical Keyboard & 12\\
102 & Optical Mouse & 30\\
103 & USB-C Cable & 100\\\bottomrule
\end{tabular}
\end{table}

An \code{INTEGER} column stores whole numbers. A \code{VARCHAR(100)} column
stores text with a maximum encoded byte length of 100 in this implementation.
A \code{BOOLEAN} holds a truth value. \code{NULL} records that a value is
missing or unknown; it is not the number zero, an empty string, or a false
Boolean. Types and nullability describe the allowed values; constraints
describe additional rules such as nonnegative stock.

\begin{lstlisting}[language=SQL]
CREATE TABLE first_product (
    product_id INTEGER PRIMARY KEY,
    name VARCHAR(100) NOT NULL,
    stock INTEGER NOT NULL CHECK (stock >= 0)
);
\end{lstlisting}

This definition creates a structure, not a row. \code{INSERT} adds rows.
\code{SELECT} retrieves values. \code{UPDATE} changes matching rows, and
\code{DELETE} removes matching rows. A table is not an ordered list: if a
result must be sorted, request \code{ORDER BY}.

\section{Keys and relationships}
\index{primary key}\index{foreign key}\index{normalization}
A \emph{primary key} identifies one row without ambiguity. It can be one
column or a combination of columns. A \emph{foreign key} names a relationship
to a key in another table. The shop stores the customer's identifier in an
order, so the order can be joined to the customer's current name without
duplicating the entire customer record.

A product can occur in several orders, and an order can contain several
products. That many-to-many relationship needs a separate table,
\code{order\_items}. Its pair \code{(order\_id, product\_id)} identifies a
product line within an order. Quantity belongs to the line, not to the
customer or product alone. Price at purchase also belongs to the line:
changing a product's current price should not rewrite past purchases.

This separation is a practical beginning to normalization: place a fact in
the table whose key determines it. Do not store a comma-separated list of
product identifiers in an order column. Separate rows make joins,
constraints, aggregation and updates possible without reparsing a text field.

\warning{Keep referenced keys stable}{At this revision MiniSQL checks normal
child inserts against referenced keys and rejects certain parent deletes,
but its foreign-key update and cascade behavior is incomplete. A parent-key
update can leave an orphan. The examples use immutable identifiers and
explicit, correctly ordered child/parent changes. Read the constraints
chapter before designing cascading workflows.}

\section{What SQL statements do}
\index{SQL}
SQL is a language for describing data structures and desired results. The
following families are useful vocabulary, not separate programs.

\begin{table}[htbp]\centering
\caption{The main kinds of work expressed in MiniSQL SQL}
\begin{tabularx}{\linewidth}{p{.18\linewidth}p{.31\linewidth}X}\toprule
Family & Examples & Purpose\\\midrule
Definition (DDL) & \code{CREATE TABLE}, \code{ALTER TABLE}, \code{DROP VIEW} & Define or change objects and rules.\\
Data changes (DML) & \code{INSERT}, \code{UPDATE}, \code{DELETE}, \code{MERGE} & Add, revise or remove rows.\\
Queries & \code{SELECT}, \code{WITH}, joins & Derive a result from stored or calculated data.\\
Transactions & \code{BEGIN}, \code{COMMIT}, \code{ROLLBACK} & Choose a boundary for related changes.\\
Access control (DCL) & \code{CREATE USER}, \code{GRANT}, \code{REVOKE} & Define principals and permissions.\\
Operations & \code{ANALYZE}, \code{SHOW STATUS}, \code{SHUTDOWN} & Inspect, maintain and manage the engine.\\\bottomrule
\end{tabularx}
\end{table}

SQL differs from an imperative row-by-row loop. ``Return the products whose
stock is less than 15'' leaves the engine free to scan a table or use an
index. You state the condition; the planner chooses an executable plan.
Small examples can look immediate, but a query still has to be parsed, bound
to real objects, checked for types and permissions, and executed.

\section{Clients, the server and the database directory}
\index{client/server}\index{session}
The \code{minisqld} program opens a database directory and listens for
connections. The \code{minisql} console connects to that server. A Python
application and the Workbench are also clients. A \emph{session} is the
server-side state of one connection: its identity, prepared plans,
transaction and other request state. Two separate \code{--query} invocations
open two separate sessions. A transaction cannot begin in one invocation
and continue in another.

\begin{figure}[htbp]
\centering
\begin{tikzpicture}[node distance=7mm and 6mm]
\node[box,text width=36mm] (shell) {Native console\\and SQL scripts};
\node[box,text width=36mm,right=of shell] (apps) {Python DB-API\\Java JDBC};
\node[box,text width=36mm,right=of apps] (gui) {Windows\\Workbench};
\node[strong,text width=120mm,below=10mm of apps] (server) {\textbf{minisqld}\quad Session, authentication, transport and resource limits};
\foreach \n in {shell,apps,gui}{\draw[arrow] (\n.south) -- (server.north -| \n.south);}
\node[box,text width=36mm,below=10mm of server.south west,anchor=north west] (sql) {Parse and bind\\SQL / catalog};
\node[box,text width=36mm,right=of sql] (plan) {Choose a plan\\Statistics / costs};
\node[box,text width=36mm,right=of plan] (execute) {Execute\\Read gate / writer};
\draw[arrow] (server.south) -- ++(0,-3mm) -| (sql.north);
\draw[arrow] (sql) -- (plan);\draw[arrow] (plan) -- (execute);
\node[strong,text width=55mm,below=12mm of sql.south west,anchor=north west] (wal) {Transactions / WAL\\Commit, recovery, checkpoints};
\node[strong,text width=55mm,right=7mm of wal] (files) {Paged files\\Tables, indexes, overflow, catalog};
\draw[arrow] (execute.south) -- ++(0,-4mm) -| (files.north);
\draw[arrow] (execute.south) -- ++(0,-4mm) -| (wal.north);
\draw[arrow] (wal) -- (files);
\end{tikzpicture}
\caption{MiniSQL's main components. Every client uses the server's SQL engine; the Workbench adds a user interface. Offline tools open a stopped database directly.}
\label{fig:architecture}
\end{figure}


The server is pointed at an already-created database directory. It does not
offer a SQL \code{CREATE DATABASE} or \code{USE database} workflow. Create a
database with \code{minisqld --init}; then start the server using the exact
directory printed by that command. Client connection paths such as
\code{/main} are protocol labels, not a way to switch the running server to
another disk database.

The directory contains authoritative catalog, table, index and WAL files as
well as metadata and derived acceleration files. Treat the database as a
whole. Copying an individual table file is not a supported way to export a
table, and renaming storage files is not the same as renaming SQL objects.
Offline tools take locks and expect the appropriate stopped-database state;
read their operation-specific instructions before running them.

\section{Durability, transactions and what an acknowledgement means}
\index{durability}\index{write-ahead log}
In autocommit mode, each successful data-changing statement has its own
transaction boundary. With an explicit transaction, you can perform several
changes and then \code{COMMIT} or \code{ROLLBACK}. MiniSQL uses a write-ahead
log (WAL), checksummed pages and recovery procedures. Committed work is
recovered using the engine's durable records; an arbitrary file copy made
while pages are changing is not a substitute for a backup procedure.

There are two different questions after a connection fails: did the server
commit the transaction, and did the application receive the reply? The
second can be false even when the first is true. An application that blindly
retries an order insertion can therefore create duplicate work. Use a stable
application identifier or unique key and a reconciliation query so you can
determine whether the operation already succeeded. This is an application
design rule, not a guarantee that a client can infer success from a lost
connection.

\section{A method for learning and diagnosing}
For each new feature, work through four questions: What result do I intend?
Which rows and keys represent it? Which MiniSQL syntax expresses it? How can
I verify the result? First run a \code{SELECT} that identifies the intended
rows. Then make the change inside a disposable transaction. Query the result
and either commit intentionally or roll back. Once correctness is established,
measure the plan and performance.

When an error occurs, preserve the smallest statement that reproduces it,
the relevant table definition, the source revision and the full error code
and text. Separate a parsing error from a constraint rejection, a connection
problem, an unsupported feature and damaged storage. The troubleshooting
chapter follows this distinction.

\tip{Checkpoint before continuing}{You should now be able to explain why
orders store a customer identifier, why order items need their own table,
why \code{NULL} differs from zero, and why a console connection is not the
database itself. The next chapter installs the programs; the guided tutorial
turns these ideas into a reproducible dataset.}
\sourcecode{src/apps/minisqld/main.ml}{Database initialization and the server's public modes.}
\sourcecode{src/minisql/client/console.ml}{A shell or script runs statements through one client session.}
\sourcecode{src/minisql/executor/dml.ml}{DML constraints and relationship validation.}
